\documentclass[UTF8,zihao=-4,oneside,openany,fontset=fandol]{ctexbook}
\usepackage[a4paper,left=30mm,right=25mm,top=27mm,bottom=27mm,headheight=16pt]{geometry}
\usepackage{amsmath,amssymb,amsthm,bm,mathtools}
\usepackage{booktabs,longtable,tabularx,array,multirow}
\usepackage{enumitem,setspace,fancyhdr,graphicx}
\usepackage{thesispalette}
\usepackage{tikz}
\usetikzlibrary{arrows.meta,positioning}
\usepackage[backend=biber,style=gb7714-2015,sorting=none,doi=true,url=false,isbn=false]{biblatex}
\addbibresource{references.bib}
\usepackage[unicode,hidelinks]{hyperref}
\usepackage{bookmark}
\setstretch{1.35}
\setlength{\parindent}{2em}
\setlength{\parskip}{0.2em}
\setlength{\emergencystretch}{3em}
\setlist{nosep,leftmargin=2em}
\ctexset{chapter={format=\centering\heiti\zihao{3},beforeskip=10pt,afterskip=22pt},section={format=\heiti\zihao{4}},subsection={format=\heiti\zihao{-4}}}
\pagestyle{fancy}\fancyhf{}
\fancyhead[C]{\small 东华大学金融硕士学位论文·方法与实证设计初稿}
\fancyfoot[C]{\thepage}
\renewcommand{\headrulewidth}{0.4pt}
\fancypagestyle{plain}{\fancyhf{}\fancyfoot[C]{\thepage}\renewcommand{\headrulewidth}{0pt}}
\numberwithin{equation}{chapter}
\newtheorem{proposition}{命题}[chapter]
\renewcommand{\proofname}{证明}
\newcommand{\RV}{\operatorname{RV}}
\newcommand{\RQ}{\operatorname{RQ}}
\newcommand{\CoVaR}{\operatorname{CoVaR}}
\newcommand{\VaR}{\operatorname{VaR}}
\newcommand{\E}{\mathbb E}
\newcommand{\Prob}{\mathbb P}
\newcommand{\F}{\mathcal F}
\newcommand{\sigmoid}{\operatorname{sigmoid}}
\newcommand{\softplus}{\operatorname{softplus}}
\newcommand{\relu}{\operatorname{ReLU}}
\newcommand{\Pending}[1]{\textbf{〔待实证填入：#1〕}}
\newcommand{\code}[1]{\texttt{\detokenize{#1}}}
\newcolumntype{Y}{>{\raggedright\arraybackslash}X}
\begin{document}
\begin{titlepage}
\centering
\vspace*{12mm}
{\heiti\zihao{1} 东华大学\par}
\vspace{14mm}
{\heiti\zihao{2} 金融硕士学位论文\par}
\vspace{19mm}
{\heiti\zihao{2}\setstretch{1.5} 基于动态 Copula-CoVaR 溢出网络\\与边特征图神经网络的中国股票市场\\已实现波动率预测研究\par}
\vspace{13mm}
{\zihao{4} 方法与实证设计初稿\quad v0.2\par}
\vfill
\begin{tabular}{rl}
姓\quad 名：& 吴立文\\[7pt]
学\quad 号：& 2251435\\[7pt]
专业方向：& 金融（MF）\\[7pt]
指导教师：& \underline{\hspace{45mm}}\\[7pt]
版本日期：& 2026 年 10 月 15 日
\end{tabular}
\vspace{15mm}
\end{titlepage}
\frontmatter
\chapter*{初稿说明}\addcontentsline{toc}{chapter}{初稿说明}
本稿依据《吴立文—开题报告》确定的题目、研究范围与六章结构展开，保留“高频测度—动态风险网络—边特征预测—样本外评价”的研究路线。第一至第四章提供可继续修订的正文、模型公式与实证设计；第五、第六章保留论证结构和待填结果位置，不以虚构数据或预设结论代替市场实验。

截至本版，已读取并记录原始数据位置索引与五份数据字典的文件哈希，并通过授权只读探测取得部分实际字段结构；字典报告版本为\code{8a69cf1}。这些证据确认了原始文档与部分数据源的schema，严格公司行动等复杂视图的核查尚未全部完成。真实小样、约三十只股票的共同八至十年覆盖及历史字段的当时可见性仍未通过验收。具体样本、日期、质量统计、预测损失和检验结果均须在真实验收与实验之后回填。

本稿将开题中尚未展开的公式具体化，并对条件方向与因果方向、标准 t-Copula 的尾部对称性、边际尺度造成的表观方向，以及风险汇总与逐边信息之间的关系作出明确区分。这些内容属于本次方法设计与推导，不应描述为开题已经给出的完整推导或已经观察到的市场发现。

版式采用中文学位论文常见的封面、摘要、目录、六章正文、参考文献和附录结构；学校正式终稿模板及细则尚未提供，故本稿不是终稿格式认证版。内部数据连接地址、配置文件、账号与源库运维细节置于独立研究任务包，不进入公开论文正文。正式提交前需由作者核对全文、数据披露授权、参考文献及学校要求。

\chapter*{摘\quad 要}\addcontentsline{toc}{chapter}{摘要}
已实现波动率预测需要同时处理资产自身波动惯性与股票之间的横截面关联。异质自回归模型通过日、周、月三个时间尺度刻画波动率动态，图预测方法则进一步引入邻居股票的信息。然而，在加入风险网络之后，预测改善究竟来自共同市场信息、总体风险强度，还是来自具体股票之间保留下来的关系，仍需要通过信息匹配的对照加以区分。

本文以中国股票市场为研究对象，构建动态 Copula-CoVaR 溢出网络与边特征图神经网络相结合的已实现波动率预测框架。高频价格用于构造五分钟采样的日度已实现方差及 HAR 特征；日收益通过 AR-GARCH-t 边际模型和二元动态 t-Copula 转化为条件尾部风险关系。模型在保留自身波动率分支的基础上，以逐边风险特征调节邻居消息，使预测器能够同时利用风险关系的强度、来源和时间变化。

围绕这一框架，本文明确点条件与事件条件 CoVaR 的不同定义，分析对称设定下上下行风险通道的镜像关系，并分离边际波动尺度与标准化依赖关系对有向边的影响。实证设计以未来一、五、二十二个交易日的日均已实现方差为目标，设置传统 HAR/HARQ、公共信息增强模型、风险汇总模型、动态相关图和底层参数边模型等对照，在相同信息集、训练损失和评估样本上检验逐边条件化的增量价值。

研究采用按时间推进的训练与预测协议，区分预测起点、标签成熟日期和数据版本；以日期级损失差、区块置信区间和模型比较检验评价结果。本文的研究重点不是预先假定复杂网络必然优越，而是建立从风险定义、关系表示到样本外证据的完整论证。数据覆盖与市场实证尚待验收完成，预测效果和适用范围将在正式结果产生后补充。

\noindent\textbf{关键词：}已实现波动率；Copula-CoVaR；风险溢出网络；边特征图神经网络；HAR；样本外预测

\chapter*{Abstract}\addcontentsline{toc}{chapter}{Abstract}
Forecasting realized volatility requires distinguishing an asset's own volatility persistence from cross-sectional information. Heterogeneous autoregressive models capture multiple time scales, while graph-based models incorporate neighboring assets. However, improvements attributed to a risk network may instead arise from common market conditions, marginal risk scales, aggregate spillover measures, or the training objective.

This study develops a forecasting framework that combines a dynamic Copula-CoVaR network with an edge-aware graph neural network for Chinese equities. Intraday prices are used to construct daily realized variance from five-minute returns. Daily returns are modeled using AR-GARCH-t margins and pairwise dynamic t-copulas. Conditional risk features then modulate messages from specific neighboring stocks, alongside a branch that preserves the target asset's own historical information.

The methodological specification distinguishes point-conditioned CoVaR from tail-event conditioning. It derives the mirror relationship between upside and downside increments under symmetric assumptions and separates marginal scale effects from standardized dependence. The empirical protocol targets the average daily realized variance over the next one, five, and twenty-two trading days. Comparisons include HAR and HARQ models, common-information controls, aggregate-risk predictors, dynamic-correlation graphs, and models receiving the underlying copula and marginal parameters.

Chronological training, label maturity, historical graph archives, and matched evaluation domains are specified before formal testing. Date-level paired losses and block-based inference are used to evaluate incremental predictive value. This draft presents the model and empirical design; live data coverage and market experiments remain to be verified. No empirical superiority is asserted before these experiments are completed.

\noindent\textbf{Keywords:} realized volatility; Copula-CoVaR; risk network; edge-aware graph neural network; HAR; out-of-sample forecasting

\tableofcontents
\chapter*{主要符号说明}\addcontentsline{toc}{chapter}{主要符号说明}
\begin{longtable}{p{0.24\textwidth}p{0.69\textwidth}}
\toprule 符号 & 含义\\\midrule\endhead
$i,j,t,h$ & 目标股票、来源股票、交易日起点、预测期限\\
$P_{i,t,s,k}$ & 股票$i$在交易日$t$、会话$s$第$k$个采样边界的价格\\
$r_{i,t,s,k},R_{i,t}$ & 会话内高频对数收益、风险模型所用日度收益\\
$\RV_{i,t},\RQ_{i,t}$ & 日度已实现方差、规则采样已实现四次变差\\
$y^{(h)}_{i,t},\widehat y^{(h)}_{i,t}$ & 未来$h$个交易日日均方差及其预测\\
$\F_t$ & 在预测起点$t$按约定可使用的信息集\\
$\sigma_{i,t+1|t},\nu_i$ & 下一期边际条件波动尺度、边际Student-t自由度\\
$\rho_{ij,t},\nu_{c,ij}$ & 起点$t$预测的二元copula相关参数、copula自由度\\
$C_{ij,t}(u,v)$ & 第一坐标为来源$j$、第二坐标为目标$i$的条件copula\\
$q_{i\leftarrow j,t}(\beta\mid\alpha)$ & 来源分位位置为$\alpha$时目标的$\beta$条件分位数\\
$\Delta^-_{i\leftarrow j,t},s_{i\leftarrow j,t}$ & 下行增量CoVaR、有符号的下行放大强度\\
$\bar s_{i\leftarrow j,t}$ & 以目标条件波动尺度标准化的下行强度\\
$A_t,\mathcal N_i(t),e_{j\to i,t}$ & 邻接矩阵、目标入邻居集合、逐边特征\\
$a_{i,t},z_{i,t},m_{i,t}$ & 风险汇总特征、节点输入、聚合消息\\
$c,s_{i,c}$ & 当次训练截止日、仅由当次训练样本估计的数值尺度\\
$\ell_{m,t},d_t$ & 模型$m$的日期级损失、基线减主模型的损失差\\
\bottomrule
\end{longtable}
\mainmatter
\chapter{绪论}
\section{研究背景与意义}
\subsection{研究背景}
金融风险预测不仅涉及资产价格变化的方向，还涉及未来价格变动的幅度及其不确定性。相同的预期收益可以对应不同的风险水平，因而波动率预测是理解和管理风险的一项基础工作。高频价格所提供的日内变动信息，使研究者能够将一个交易日内部的价格路径转化为日度已实现测度，并在较低频率上开展预测研究。已实现测度的构造、条件波动的建模与预测分布之间的联系，构成了高频金融计量的重要研究基础\cite{andersen2003}。

已实现波动率的预测通常从资产自身历史出发。HAR 模型以日、周、月尺度的滞后波动率构造简约的预测关系，既便于估计，又为加入其他信息提供了清楚的参照\cite{corsi2009}。在多资产环境中，研究问题进一步扩展为：在某只股票的自身波动历史之外，其他股票当下的风险状态是否还包含有助于预测的信息？回答这一问题并不意味着应当舍弃自身历史，而是需要检验新增横截面信息能否在强基线之上提供增量。

横截面信息可以通过不同方式进入预测模型。市场或行业平均波动率提供总体状态，风险输入与输出指标概括一个节点在网络中的位置，逐边模型则保留不同邻居与目标股票之间的对应关系。三种表示保留的信息不同，复杂度也不同。本文关心的不是抽象意义上“使用网络是否更先进”，而是在风险已经可以被汇总为少数变量时，进一步保留每一条关系是否值得。

现有研究已经将图神经网络用于已实现波动率预测，也已将动态 Copula-CoVaR 生成的风险汇总指标引入 HAR 框架\cite{zhang2025,wang2026}。因此，本文不将“HAR 与 GNN 相结合”或“CoVaR 进入波动率预测”本身视为全新问题，而将研究重点放在关系信息的压缩与保留：当风险总量相近时，来自不同股票、具有不同时间变化的条件风险关系，是否应以不同方式参与目标股票的预测？

\subsection{理论意义}
本文的理论意义首先在于将风险度量与预测任务之间的联系具体化。CoVaR 是一个条件风险量，并不直接等于未来已实现方差。本文将其视为基于日收益条件分布得到的关系特征，检验其能否帮助预测另一个明确界定的对象，而不将“度量具有经济含义”作为“必然具有预测能力”的替代证据。

其次，本文区分关系表示的不同层次。二值图描述哪些股票相连，标量边描述关系强度，多维边进一步记录标准化强度及其变化。通过保持节点信息与骨干模型一致，可以将这种表示差异转化为可比较的模型限制，而不是一次性改变图、损失函数和模型容量之后，再笼统归因于风险网络。

最后，本文明确讨论风险语义与模型结构的边界。条件方向不等于因果方向；动态相关不等于上下尾不对称；未经标准化的双向 CoVaR 差异也可能主要体现目标股票波动尺度的不同。对这些概念作出区分，可以使研究结论与模型实际识别的对象保持一致。

\subsection{实践意义}
本文预期为风险预测提供一条可审计的实现路径。对于使用预测结果的研究人员，除了一个最终方差预测值，还需要知道输入数据的时段范围、风险边的估计截止日、图估计失败时的处理方式，以及模型在不同市场状态下是否稳定。将这些信息纳入实验记录，有助于区分模型误差、数据问题与使用时点问题。

实践价值需要建立在相应证据之上。预测误差的降低并不自动意味着交易收益提高，也不自动构成 VaR 或 ES 的有效性证明。本文首先关注预测任务本身；资产配置、波动率择时与风险管理应用只有在收益口径、信号可用时间和交易约束都可核验的情况下，才作为后续扩展分析。

\section{研究内容与研究问题}
\subsection{研究内容}
本文沿开题报告确定的研究路线\cite{proposal2026}，将工作组织为数据测量、风险关系估计、预测模型与样本外比较四个环节。第一，以公司分钟价格序列构造五分钟采样的日度已实现方差，明确集合竞价、午休、停牌、缺失和收益调整的处理。第二，用相同股票的日收益估计边际分布和动态依赖，从同一二元模型计算两个条件方向的风险量。第三，在自身 HAR 分支之外设置浅层逐边消息聚合，同时建立使用相同风险汇总信息的竞争模型。第四，按时间推进完成不同期限和训练损失的预测，结合关系破坏、尺度标准化和底层参数对照分析增量来源。

\subsection{研究问题}
本文围绕以下四个相互衔接的问题展开。第一，条件尾部风险关系在原始尺度与标准化尺度下呈现何种动态差异，其中多少可以由边际波动尺度或分布形状解释？第二，相比公共市场信息和风险输入、输出汇总，逐边关系是否提供额外的样本外预测价值？第三，在控制动态相关、边际风险与模型容量后，CoVaR 表征以及多维边条件化分别起什么作用？第四，上述价值是否随预测期限、市场状态、数据采样和估计窗口发生变化？

这里的“方向性”首先指条件设定中来源与目标的区分，不预设已识别真实市场中的单向传染。对问题的回答将依赖预先设定的比较和实际数据，而不依赖网络图的视觉复杂度。

\section{研究方法与技术路线}
\subsection{研究方法}
文献研究用于确定研究与已有工作的关系，尤其是图上的非线性溢出预测和 CoVaR 风险汇总预测之间的交叉。数学分析用于澄清条件分位数、对称性与尺度分解，使实现能够接受可复核的单元测试。实证比较则将不同信息集合与函数形式放入统一的时间外评价协议，分别检验其作用。

技术实现采用两条数据路径。分钟价格进入日度已实现方差及 HAR 特征，日收益进入动态条件风险边。两条路径在预测起点汇合，但未来真实收益与未来已实现方差不进入任何输入。风险图参数按预设频率重估，状态在可用观测到达后过滤，预测模型只用已经成熟的标签更新。

\subsection{论证顺序}
论文的论证从最简自身基线开始，依次增加公共信息、风险汇总、具体关系和多维边输入。这个顺序使每一步新增内容都有相应对照。若逐边模型只能优于未加入公共信息的 HAR，结论就限于横截面信息有效；只有进一步超过等信息风险汇总与动态相关等控制，才支持更强的关系级表征主张。

\section{研究设计的主要特点}
本文拟形成三个层面的研究贡献。其一，围绕“风险汇总是否压缩了对预测有用的关系”构建模型与实验，使研究动机直接对应可验证的问题。其二，将条件风险边的定义、标准化与归因控制纳入方法设计，明确风险尺度、相关参数和边表示各自的作用。其三，在公司高频数据场景中建立可追溯的数据与时间协议，使模型比较能够对应相同的样本、标签和信息集。

上述内容是本文的设计特点。对预测优越性、跨状态稳定性和应用价值的结论仍须由第五章的真实实证支持，不将方法组合或工作量本身视为已经证实的效果。

\section{论文结构}
全文分为六章。第一章提出研究背景、问题与总体路线。第二章梳理已实现波动率预测、动态条件风险及图预测文献。第三章给出研究对象、风险模型、边表示、预测表达式和必要数学性质。第四章说明公司数据合同、样本构造与实验协议。第五章按整体表现、贡献来源和稳健性组织实证结果。第六章根据证据总结结论并讨论实践启示与适用边界。

\chapter{文献综述}
\section{已实现波动率预测研究}
\subsection{低频条件波动与高频已实现测度}
GARCH 模型通过收益创新和滞后条件方差描述波动动态，为低频收益的条件异方差建模提供了基本工具\cite{bollerslev1986}。高频已实现测度则从日内价格路径出发，将多个高频收益的平方累计为日度风险代理。两者并非互相排斥：条件波动模型刻画日收益分布，已实现测度用于构造和预测日内变化；本文分别将其用于风险边与预测目标。

在适当的价格过程及观测条件下，已实现方差与二次变差之间存在理论联系\cite{andersen2003}。实际数据仍涉及有限采样、价格噪声、停牌与不规则时段等问题。因此，高频数据行数更多并不自动意味着日度测度更准确；一个明确定义、来源一致的五分钟价格路径，通常比将不同时段和来源的记录机械相加更适合可复核比较。本文选择五分钟作为起始采样尺度，并以替代采样频率检验结果对测度的敏感性，不将其预设为所有数据场景的最优选择。

\subsection{HAR、测量误差与自身信息扩展}
HAR 将不同时间尺度的滞后已实现波动率放入简约回归，以近似表示波动率较长的持续性\cite{corsi2009}。其价值不仅体现在基准预测能力，还体现在结构清楚：任何新增公共因子、图聚合项或非线性分支，都可以与自身历史形成直接对照。

Bollerslev、Patton 和 Quaedvlieg 将已实现测度的误差大小纳入预测系数调整，提出 HARQ 等模型\cite{bpq2016}。本文采用其日度四次变差交互项的简约版本作为计量基线。HARQ 中的 Q 指四次变差相关的误差调整，并非 QLIKE 训练的缩写；后者在本文单独标注为训练损失。

对于自身信息，半方差或收益符号能够提供不同于总方差的描述。本文设置相应的自身特征扩展对照，是为了避免把本来可以由目标股票自身状态解释的改善全部归因于外部风险边。本研究不将所有自身特征扩展同时堆入主模型，而按照信息匹配原则在明确的比较组中使用。

\subsection{预测损失与模型比较}
方差通常通过代理量评价，损失函数对比较结论具有重要影响。Patton 讨论了使用不完美波动率代理时的预测比较条件\cite{patton2011}。这一研究为考察不同损失提供了依据，但不意味着任意来源、任意处理过的 RV 都自动满足相关条件。

在本文中，MSE 强调原始方差单位下的平方偏差，QLIKE 依赖实际与预测的比值。两者既是可能的训练准则，也是需要明确区分的评价准则。模型使用 QLIKE 训练之后，仍需与使用同一准则训练的无图模型比较。由此获得的差值，才更接近网络表征的增量，而不是训练方式改变的混合结果。

\section{波动率溢出与动态 Copula-CoVaR 研究}
\subsection{条件风险与依赖结构}
多资产风险建模需要区分单个资产的边际分布与资产之间的依赖关系。动态相关模型为时间变化的关联提供简约参数化\cite{engle2002}；Copula 表示则将边际分布与联合依赖分开，使不同边际下的条件风险计算可以在统一的概率坐标中进行。

Adrian 和 Brunnermeier 提出的 CoVaR 从某一主体处于特定状态的条件出发，考察另一主体或系统的风险分位数，增量形式则比较困境与基准状态\cite{adrian2016}。本文将这一思路用于股票对之间的条件风险关系，但研究对象和目标不同于原有系统性风险度量：股票对风险边在这里是未来已实现方差的预测输入，而不是对宏观系统重要性的直接认定。

\subsection{条件定义与方向语义}
来源收益恰好处于一个分位数位置，与来源收益落入该分位数以下的整个尾部事件，是两个不同条件。Mainik 和 Schaanning 讨论了不同条件设定下 CoVaR 对依赖关系的一致性问题\cite{mainik2012}。本文因此将点条件与事件条件分开定义、保存和比较，不在程序中只用一个 CoVaR 名称掩盖二者差异。

条件方向也不同于时间领先或因果方向。$i\leftarrow j$ 表示在$j$的给定分位状态下计算$i$的条件分位数；即使同时得到两个方向数值不同，也还需要分析这种差异是否由边际尺度造成。本文所称“风险溢出网络”沿用开题中的经济解释，但对可识别内容作出上述限制。

\subsection{尾部对称性与非对称扩展}
标准 t-Copula 具有上下尾对称的结构\cite{demarta2005}。在对称边际和镜像分位条件下，仅将分位参数从下尾改到上尾并不会产生独立的第二种风险机制。本文把这一性质写入模型定义与前置测试，以防将数值误差或重复变量解释成市场非对称证据。

允许上下尾不同的 Copula 可以用于进一步研究，例如 Patton 对非对称依赖的建模\cite{patton2006}。本文将此类模型作为明确命名的扩展，而不是在原型效果不理想后无记录地替换。上下尾不对称本身也不意味着来源与目标之间的非交换性，两者需分别检验。

\subsection{风险汇总与已实现波动率预测}
与本文最直接相邻的工作是 Wang、Yuan 和 Qu 的极端市场状态切换模型。公开摘要与引言显示，该研究基于动态 Copula-CoVaR 构建上、下行风险输出和接收指标，并将其与 HAR 相结合\cite{wang2026}。这一证据已经说明，CoVaR 与波动率预测的结合并非本文首次提出。

本文的区别在于进一步询问：当节点输入、输出的风险总量已被模型使用后，具体邻居与边特征之间的对应关系是否仍有价值。由于目前未取得该工作的完整方法公式与实现，本文先构造定义透明的风险汇总基线，不将其命名为该文的严格复现；取得全文后再通过公式对应表补充复现版本。

\section{图神经网络与金融图预测}
\subsection{图结构与边输入}
图模型把股票作为节点，通过邻居聚合描述横截面联系。GLASSO 在精度矩阵上进行稀疏估计，是统计构图的一种方式\cite{friedman2008}。其所反映的是特定模型假设下的条件关联结构，不应被直接解释为因果关系。本文保留其作为竞争构图方式，而不是先验否定其风险描述能力。

消息传递框架允许模型根据节点和边的信息构造消息并聚合到目标节点\cite{gilmer2017}。因此，边特征机制本身不是本文发明的通用图学习算子。本文的方法贡献在于为这一任务选择明确的风险边表示，并在统一实验中检验其是否比标量权重或风险汇总更有用。

\subsection{GNNHAR 与训练准则}
Zhang 等将图上的非线性溢出纳入多资产已实现波动率预测，并同时研究邻居范围和训练准则\cite{zhang2025}。这一工作提示本文应把信息集合、函数形式和训练损失分别控制，而不把它们的联合改变压缩成单一的“GNN 改进”。

本文设置两组相关对照：一组尽可能对应原文的图与模型设定，另一组在统一输出约束、节点变量和训练预算下替换构图及边表示。前者用于说明与既有文献的联系，后者用于更精确地识别本项目内部的增量来源。预印本与最终发表版的实现差异亦需在复现说明中记录。

\section{文献评述与本文定位}
\begin{table}[htbp]\centering\small
\caption{相关研究与本文研究问题的对应关系}
\begin{tabularx}{\textwidth}{p{0.2\textwidth}YY}
\toprule 研究线索 & 已提供的基础 & 本文进一步检验的内容\\\midrule
HAR 与 HARQ & 多尺度自身历史、测量误差调整 & 新增横截面信息能否超过强自身基线\\
GNNHAR & 图上的非线性聚合及损失比较 & 风险图与多维边是否有独立增量\\
CoVaR 风险汇总预测 & 风险输出/接收及状态信息 & 同风险汇总下保留具体关系的价值\\
Copula 风险度量 & 边际与依赖分离、条件分位数 & 风险语义、尺度与模型识别边界\\
边特征消息传递 & 关系依赖的消息函数 & 对这一金融任务的有限样本适用性\\\bottomrule
\end{tabularx}
\end{table}
本文不以“已有方法都忽略尾部风险”或“相关图无法识别任何风险”作为出发点。更具体的定位是：在既有预测与风险度量基础上，建立从风险汇总到逐边条件化的可检验联系。若新增表示能够稳定提高预测能力，说明这种关系信息在所研究样本与协议下值得保留；若风险汇总已足够，则说明更简约的信息压缩在相应条件下具有竞争力。两类结果都应通过相同实验标准判断。

\chapter{理论基础与模型构建}
\section{预测任务、信息集与已实现测度}
\subsection{预测对象与信息边界}
设研究股票集合为$\mathcal V$，其规模为$N$；$t$表示共同交易日历上的一个预测起点，而不是某只股票的第$t$条有效记录。用$\F_t$表示在日$t$信息完成后、按照数据与版本协议允许使用的信息集。本文预测未来$h$个交易日的日均已实现方差：
\begin{equation}\label{eq:target}
y_{i,t}^{(h)}=\frac1h\sum_{\ell=1}^{h}\RV_{i,t+\ell},\qquad h\in\{1,5,22\}.
\end{equation}
三个期限分别使用直接预测模型。$h=5$和$h=22$并非只预测第5天或第22天，也不是将未来收益率相加后求平方。若目标窗口内任何一天的测度不可用，则该股票该起点的完整标签不可用，不改成剩余有效日的平均值。

本文沿用开题中的“已实现波动率”名称，但$\RV$在数学上是收益率平方和，即方差量。其平方根才具有与收益标准差相同的量纲。除明确说明的展示指标外，预测目标不年化，所有收益均使用小数单位。

\subsection{会话内已实现方差与四次变差}
将每个交易日分为经认证的上午和下午交易会话，记为$s\in\{\mathrm{AM},\mathrm{PM}\}$。设$P_{i,t,s,0}$为可追溯的会话起始锚点，$P_{i,t,s,k}$为第$k$个五分钟采样边界的价格。高频对数收益和主测度定义为
\begin{align}
 r_{i,t,s,k}&=\log P_{i,t,s,k}-\log P_{i,t,s,k-1},\label{eq:intradayreturn}\\
 \RV_{i,t}&=\sum_s\sum_{k=1}^{M_s}r_{i,t,s,k}^{2}.\label{eq:rv}
\end{align}
主测度不额外包含隔夜和午间重开跳变；收盘集合竞价是否在最后边界内，由数据合同明确。对于完整、规则的会话内采样集合，记$M=\sum_s M_s$，定义
\begin{equation}\label{eq:rq}
\RQ_{i,t}=\frac M3\sum_s\sum_{k=1}^{M_s}r_{i,t,s,k}^{4}.
\end{equation}
式\eqref{eq:rq}沿用规则采样四次变差的形式\cite{bpq2016}。若大量区间缺失、起始锚点不规则或额外并入隔夜长间隔，不应只调整$M$就假定保留相同的测量误差解释。本文将此类情况单独标记或使用独立扩展测度。

为检验时段选择，另保存午间收益$r^{L}_{i,t}$和经调整口径确认的隔夜收益$r^{O}_{i,t}$。扩展测度可分别写为$\RV_{i,t}+(r^L_{i,t})^2$及进一步加入$(r^O_{i,t})^2$。这些变量具有不同的覆盖范围，在实验注册中使用不同名称，不在同一目标列混用。

\subsection{HAR 特征与基准模型}
以共同交易日历定义三个尺度的历史特征：
\begin{equation}\label{eq:harfeatures}
x^{H}_{i,t}=\left(\RV_{i,t},\ \frac15\sum_{\ell=0}^{4}\RV_{i,t-\ell},\ \frac1{22}\sum_{\ell=0}^{21}\RV_{i,t-\ell}\right)^{\!\top}.
\end{equation}
传统线性 HAR 基准为
\begin{equation}\label{eq:har}
y_{i,t}^{(h)}=b_{i,h}+\beta_h^{\top}x^{H}_{i,t}+u_{i,t}^{(h)}.
\end{equation}
本文同时保存逐股 OLS 与具有股票截距的合并估计版本，注册时区分参数共享方式。任何模型都不能通过删除中间缺失日，将式\eqref{eq:harfeatures}变成最近5或22个“有记录日”的均值。

HARQ 的简约版本在日度项上加入四次变差交互\cite{bpq2016}：
\begin{align}
y_{i,t}^{(h)}={}&b_{i,h}+\beta_{d,h}\RV_{i,t}
 +\beta_{w,h}\RV_{i,t}^{(w)}+\beta_{m,h}\RV_{i,t}^{(m)}\nonumber\\
&+\beta_{q,h}\big(\sqrt{\RQ_{i,t}}-\overline{\sqrt{\RQ_i}}_{\mathrm{train}}\big)\RV_{i,t}
+u_{i,t}^{(h)}.\label{eq:harq}
\end{align}
其中中心化均值只由当次训练样本估计。式\eqref{eq:harq}的日度形式对应原文的简约 HARQ，较长期限采用同特征直接预测，是本文的实施扩展。HARQ 与使用 QLIKE 训练的 HAR 在模型注册中分别命名。

\section{日收益边际分布与动态依赖}
\subsection{风险输入的收益口径}
记$R_{i,t}$为风险模型所用日度收益。公司日线提供的前收字段被报告为当日交易所参考价$B_{i,t}$，不保证等于前一交易日未复权收盘价。本文的候选主输入为
\begin{equation}\label{eq:reference_return}
R^{\mathrm{ref}}_{i,t}=\log\frac{P^{\mathrm{close}}_{i,t}}{B_{i,t}},\qquad B_{i,t}>0.
\end{equation}
该输入必须经过普通日和公司行动日核对后启用。它表示相对参考价的价格变化，不直接称为投资者总回报。不同调整收益口径进入独立稳健性实验，不在同一风险历史中静默拼接。风险模型包含的日间与隔夜信息可以不同于会话内 RV 的覆盖范围；这是输入与目标的任务安排，不表示二者是同一风险量。

\subsection{AR-GARCH-t 边际模型}
对股票$i$建立AR(1)-GARCH(1,1)边际：
\begin{align}
R_{i,t}&=c_i+\phi_iR_{i,t-1}+\epsilon_{i,t},\label{eq:ar}\\
\epsilon_{i,t}&=\sigma_{i,t}\xi_{i,t},\qquad \xi_{i,t}\sim G_{\nu_i},\label{eq:innov}\\
\sigma_{i,t}^{2}&=\omega_i+a_i\epsilon_{i,t-1}^{2}+b_i\sigma_{i,t-1}^{2}.\label{eq:garch}
\end{align}
$G_{\nu}$表示单位方差Student-t分布，而$T_{\nu}$表示通常的Student-t分布函数，两者满足
\begin{equation}\label{eq:unit_t}
G_{\nu}(x)=T_{\nu}\!\left(x\sqrt{\frac{\nu}{\nu-2}}\right),\qquad
G_{\nu}^{-1}(u)=\sqrt{\frac{\nu-2}{\nu}}T_{\nu}^{-1}(u),\quad \nu>2.
\end{equation}
约束为$\omega_i>0$、$a_i,b_i\geq0$、$a_i+b_i<1$及$|\phi_i|<1$。式\eqref{eq:garch}保留标准 GARCH 的尺度动态\cite{bollerslev1986}，t边际用于本研究对厚尾形状的参数化。

在起点$t$，令$m_{i,t+1|t}=c_i+\phi_iR_{i,t}$，则下一期条件边际分布为
\begin{equation}\label{eq:forecast_margin}
F_{i,t+1|t}(x)=G_{\nu_i}\!\left(\frac{x-m_{i,t+1|t}}{\sigma_{i,t+1|t}}\right).
\end{equation}
拟合和过滤只使用$\F_t$中的观测。历史概率积分变换为$U_{i,s}=G_{\nu_i}(\xi_{i,s})$。单位方差t与普通t、边际自由度与copula自由度均分别保存，不能因软件接口名称相近而混用。

\subsection{二元 DCC 型 t-Copula}
每个无序股票对$\{i,j\}$只拟合一个联合依赖模型。记$\boldsymbol\xi_{ij,s}=(\xi_{j,s},\xi_{i,s})^\top$，以边际标准化创新驱动DCC型递推：
\begin{align}
Q_{ij,s+1}&=(1-a_{ij}-b_{ij})\bar Q_{ij}
+a_{ij}\boldsymbol\xi_{ij,s}\boldsymbol\xi_{ij,s}^{\top}+b_{ij}Q_{ij,s},\label{eq:dcc}\\
\mathcal R_{ij,s+1}&=D_{ij,s+1}^{-1/2}Q_{ij,s+1}D_{ij,s+1}^{-1/2},\label{eq:dccnorm}
\end{align}
其中$D$为$Q$的对角矩阵，$\bar Q$在当前历史估计窗口内确定且保持正定，$a_{ij},b_{ij}\geq0$、$a_{ij}+b_{ij}<1$。动态相关的基本思想来自DCC\cite{engle2002}；本研究选定的是以边际创新驱动、以copula似然拟合依赖参数的两阶段二元版本。

令$\rho_{ij,t}=[\mathcal R_{ij,t+1|t}]_{12}$，即$t$时点给出的下一期相关状态，$\nu_{c,ij}$为copula自由度。定义
\begin{equation}\label{eq:tcopula}
C_{ij,t}(u,v)=T_{\nu_{c,ij},\rho_{ij,t}}
\left(T_{\nu_{c,ij}}^{-1}(u),T_{\nu_{c,ij}}^{-1}(v)\right),
\end{equation}
其中第一坐标始终为来源$j$，第二坐标为目标$i$。依赖参数目标为有效历史样本的copula对数密度和，密度为二元t密度除以两个对应一元t密度的乘积。参数约每21个交易日重新估计，期间只用新观测过滤状态。

因各二元模型独立估计，其相关参数拼接后不一定形成合法的$N$维相关矩阵。因此，本文使用的是关系级图特征，而不是已经认证的全维联合收益分布。不在缺少进一步约束时直接用该矩阵构造最小方差组合。

\subsection{缺失观测与历史图}
停牌或数据缺口不意味着收益创新为零，也不能把缺口两侧的收益压缩成相邻交易日。本文初版采用分段过滤：在连续有效日收益片段内过滤状态，每段设置独立初值及22个观测的预热；参数目标由各片段预热后的有效项组成。二元依赖只在双方创新均有效的连续片段上递推。最近连续片段不足时，风险边标记为不可用，而不是填零。

该规则是一项保守的分段条件估计设计，不声称等价于对所有缺失收益的精确联合似然。窗口内有效项门槛与边覆盖损失在开发阶段报告，并在测试之前冻结。历史风险图按起点归档，当前重新估计的参数不得回写更早日期的图，以免使训练输入包含当时尚未出现的信息。

\section{CoVaR 定义与条件分位数计算}
\subsection{点条件 CoVaR 与增量形式}
给定$\alpha,\beta\in(0,1)$，点条件分位数定义为
\begin{equation}\label{eq:point_covar}
q_{i\leftarrow j,t}(\beta\mid\alpha)
=Q_{\beta}\!\left(R_{i,t+1}\mid U_{j,t+1}=\alpha;\F_t\right).
\end{equation}
此处条件于连续变量的一个点，通过模型的正规条件分布定义，不是直接从历史样本中寻找“收益恰好等于VaR”的记录。采用$\alpha=\beta=0.05$的候选主设定，增量为
\begin{align}
\Delta^-_{i\leftarrow j,t}&=q_{i\leftarrow j,t}(\beta\mid\alpha)
-q_{i\leftarrow j,t}(\beta\mid0.5),\label{eq:delta_down}\\
\Delta^+_{i\leftarrow j,t}&=q_{i\leftarrow j,t}(1-\beta\mid1-\alpha)
-q_{i\leftarrow j,t}(1-\beta\mid0.5).\label{eq:delta_up}
\end{align}
式\eqref{eq:delta_down}与式\eqref{eq:delta_up}比较困境与中位数点参考态，沿用CoVaR增量的思想\cite{adrian2016}。定义有符号下行放大强度$s_{i\leftarrow j,t}=-\Delta^-_{i\leftarrow j,t}$。$s>0$表示该定义下目标下尾更低，$s<0$表示相对参考态的缓释，不能统一取绝对值后均称为风险放大。

\subsection{t-Copula 的闭式计算}
省略股票对和时间下标，记$\nu_c$为copula自由度，$a_{\alpha}=T_{\nu_c}^{-1}(\alpha)$。由二元t条件分布可得
\begin{align}
z_{\beta\mid\alpha}
&=\rho a_{\alpha}+\sqrt{\frac{(\nu_c+a_{\alpha}^{2})(1-\rho^{2})}{\nu_c+1}}
T_{\nu_c+1}^{-1}(\beta),\label{eq:conditional_t}\\
v_{\beta\mid\alpha}&=T_{\nu_c}(z_{\beta\mid\alpha}),\label{eq:prob_position}\\
q_{i\leftarrow j,t}(\beta\mid\alpha)
&=m_{i,t+1|t}+\sigma_{i,t+1|t}G_{\nu_i}^{-1}(v_{\beta\mid\alpha}).\label{eq:quantile_back}
\end{align}
式\eqref{eq:conditional_t}在copula的普通t坐标上运算，式\eqref{eq:quantile_back}才返回目标的单位方差t边际。由同一组$\rho,\nu_c$和相应目标边际得到反向条件量，不重新拟合另一套二元copula来制造方向差。

这种闭式计算避免将大量蒙特卡洛模拟作为日度构图的必要步骤。其数值实现需要检查分位函数单调性、参数边界和极端概率；裁剪仅为端点数值保护，不作为人为压低尾部风险的措施。

\subsection{事件条件稳健性}
点条件与整个尾部事件是不同风险量\cite{mainik2012}。对于下尾事件$U_j\leq\alpha$，目标概率位置$v^-_{\alpha,\beta}$满足
\begin{equation}\label{eq:event_lower}
\frac{C_{ij,t}(\alpha,v^-_{\alpha,\beta})}{\alpha}=\beta.
\end{equation}
对于上尾事件$U_j\geq1-\alpha$，其$(1-\beta)$分位位置满足
\begin{equation}\label{eq:event_upper}
\frac{v^+_{\alpha,\beta}-C_{ij,t}(1-\alpha,v^+_{\alpha,\beta})}{\alpha}=1-\beta.
\end{equation}
求得$v$后经式\eqref{eq:forecast_margin}的逆变换返回收益坐标。事件版本下，下尾参考态设为$U_j\leq0.5$，上尾参考态设为$U_j\geq0.5$，并分别计算增量。本文将此版本命名为“尾部事件相对同侧半分布”，不能称为中位数点参考态。两种版本在实验和文件中分开保存。

\section{风险边的结构性质与解释边界}
\subsection{上下尾的镜像关系}
\begin{proposition}[对称设定的镜像性质]\label{prop:symmetry}
若copula具有中心对称性及连续条件密度，条件分位数唯一，目标边际连续且关于其条件位置$m_i$对称，且上下尾使用镜像概率和相同中位数点参考态，则
\begin{equation}\label{eq:symmetry}
\Delta^+_{i\leftarrow j,t}=-\Delta^-_{i\leftarrow j,t}.
\end{equation}
\end{proposition}
标准t-Copula与本文的对称t边际满足相应结构条件\cite{demarta2005}。证明见附录。该结果意味着在忠实原型中，上、下行正向表示的风险强度重复。本文保留两列用于数值审计，但不把重复列解释为两个独立的市场机制。允许尾部非对称的模型需要另行指定，例如相应的非对称copula扩展\cite{patton2006}。

\subsection{边际尺度与条件方向}
\begin{proposition}[同形边际的尺度分解]\label{prop:scale}
若二元copula可交换，两端条件边际属于同一标准化形状的位置尺度族，且使用相同分位数定义，则存在对称系数$k_{ij,t}=k_{ji,t}$，使
\begin{equation}\label{eq:scale_decomposition}
s_{i\leftarrow j,t}=\sigma_{i,t+1|t}k_{ij,t},\qquad
s_{j\leftarrow i,t}=\sigma_{j,t+1|t}k_{ij,t}.
\end{equation}
\end{proposition}
因此定义标准化强度
\begin{equation}\label{eq:standard_strength}
\bar s_{i\leftarrow j,t}=\frac{s_{i\leftarrow j,t}}{\sigma_{i,t+1|t}}.
\end{equation}
命题\ref{prop:scale}的条件下，双向$\bar s$相等。若边际自由度不同，标准化之后仍可能保留边际形状引起的差异。因此，本文同时保存原始尺度、标准化尺度与边际自由度，而不单凭原始矩阵非对称认定存在独立的方向结构。

此外，同一目标的所有原始正权重都乘有相同$\sigma_i$时，按该目标的权重和归一化会抵消这一尺度；逐目标top-$k$还可能将原本对称的权重矩阵变成非对称邻接支持。这两种变化来自处理规则，必须与估计出的条件关系区分。

\subsection{独立性与零相关}
独立copula为$C(u,v)=uv$，其点条件和事件条件目标分布均不受来源状态影响，故对应增量为零。但有限自由度t-Copula即使$\rho=0$，也不一般等于独立copula\cite{demarta2005}。程序的独立性测试因此使用显式独立模型，而不是把相关参数设为零后期待所有风险增量都消失。

上述命题是本文所选设定的代数性质与解释工具，不是关于中国市场不存在风险传导的命题，也不作为新的金融理论创新加以夸大。

\section{风险汇总与逐边表示}
\subsection{风险网络与边特征}
在起点$t$，从估计有效且下行放大强度为正的候选关系中，为每个目标股票选取至多$k$条入边。主开发候选为$k=5$，按$\bar s_{i\leftarrow j,t}$排序，平局按固定证券键处理。定义
\begin{equation}\label{eq:adjacency}
A_{ij,t}=\mathbf1\{j\in\mathcal N_i(t)\},\qquad A_{ii,t}=0,
\end{equation}
即邻接矩阵的行表示目标、列表示来源。核心边向量为
\begin{equation}\label{eq:edgefeatures}
e_{j\to i,t}=
\left(s_{i\leftarrow j,t},\ \bar s_{i\leftarrow j,t},\ \bar s_{j\leftarrow i,t},\
\bar s_{i\leftarrow j,t}-\bar s_{i\leftarrow j,t-5}\right)^{\!\top}.
\end{equation}
第五个交易日前的边必须来自当时归档版本，不用当前参数重算过去。差分项缺失时保留指示并按统一规则处理，不能将其解释为零变化。各边分量只用训练样本进行数值缩放，不以测试分位数截断。

\subsection{风险输入与输出汇总}
记$\mathcal J_{i,t}^{\mathrm{in}}$和$\mathcal J_{i,t}^{\mathrm{out}}$为所有估计有效的入向、出向关系，$[x]_+=\max(x,0)$。以有效边均值定义
\begin{align}
I_{i,t}&=\frac{1}{|\mathcal J_{i,t}^{\mathrm{in}}|}\sum_{j\in\mathcal J_{i,t}^{\mathrm{in}}}[s_{i\leftarrow j,t}]_+,\label{eq:inaggregate}\\
O_{i,t}&=\frac{1}{|\mathcal J_{i,t}^{\mathrm{out}}|}\sum_{j\in\mathcal J_{i,t}^{\mathrm{out}}}[s_{j\leftarrow i,t}]_+.\label{eq:outaggregate}
\end{align}
标准化汇总$\bar I,\bar O$以$\bar s$替换$s$。风险汇总输入$a_{i,t}$包含这四个量和入、出向有效边比例；无有效边时汇总设为缺失并提供指示，不能将“没有估计”与“估计出零风险”混淆。另设与主图相同支持下的汇总版本，以控制图稀疏化的影响。

这些是本文明确自定义的汇总基线，并非对已有EMRS公式的猜测性复现。对负向缓释关系保留原始记录，必要时作独立汇总稳健性；主放大网络不将负值绝对化。

\subsection{为什么汇总可能不足}
汇总表示压缩了边与来源股票之间的对应关系。考虑两来源节点1、2和两目标节点3、4，边权分别为$(a,b)$与$(b,a)$。同时交换两目标的两个入边权后，每个节点的总入、出权重都不变，但目标3的加权邻居状态由$ax_1+bx_2$变为$bx_1+ax_2$。二者差为$(a-b)(x_1-x_2)$。

这个构造说明风险总量并不在数学上决定加权邻居状态，因而逐边表示存在保留更多信息的可能。它不是实际股票风险网络的构造证明，更不保证新增信息在有限样本中能带来预测收益。是否值得保留仍由等信息汇总模型与逐边模型的样本外差值决定。

\section{Edge-aware CoVaR-GNN-HAR 预测模型}
\subsection{节点信息与两种模型版本}
为对应开题架构，设置最小版本，节点仅使用HAR特征。为进行严格归因，另设完整匹配版本，其节点向量$z_{i,t}$包含：HAR特征、剔除自身的样本市场与可认证行业平均RV、边际条件风险变量、同样的风险汇总$a_{i,t}$及缺失指示。若历史行业不可认证，同一比较组全部删去行业变量。

记$s_{i,c}>0$为当次训练截止$c$之前估计的股票RV均值，HAR输入以其缩放。其他变量使用训练统计量标准化。两个版本分别注册为\code{M01_BASIC}和\code{M01_FULL}；核心“超过风险汇总”的比较以完整匹配版本为准，不将不同节点信息集混成一个模型结果。

图\ref{fig:matched_architecture}展示完整模型与匹配基线的共同输入和差异。它是模型结构示意，不包含真实市场样本或预测表现；所有经验差值均留待第五章填入。
\begin{figure}[htbp]
\centering
\begin{tikzpicture}[
  box/.style={draw=ThesisText,line width=.65pt,rounded corners=2pt,
    align=center,text=ThesisText,font=\small,inner sep=7pt},
  arrow/.style={-{Stealth[length=2mm]},draw=ThesisText,line width=.8pt}]
\node[box,fill=ThesisLilac,text width=12.8cm] (common) at (0,0)
  {共同节点信息 $z_{i,t}$\\
  HAR、剔除自身的公共RV、边际风险、风险汇总与有效性标志};
\node[box,fill=ThesisMist,text width=5.5cm,minimum height=1.6cm] (own) at (-3.6,-1.8)
  {共同自身分支\\HAR线性项 $+$ 自身MLP\\两模型使用同样的输入与形式};
\node[box,fill=ThesisSalmon,text width=5.5cm,minimum height=1.6cm] (edge) at (3.6,-1.8)
  {逐边消息分支\\来源节点、目标节点与边特征\\sigmoid门控 $+$ 度数均值};
\node[box,dashed,fill=ThesisMatched,text width=5.5cm,minimum height=1.35cm] (matched) at (-3.6,-3.8)
  {\texttt{B06\_MATCHED}\\共同自身分支 $\rightarrow$ 正值预测};
\node[box,fill=ThesisFull,text width=5.5cm,minimum height=1.35cm] (full) at (3.6,-3.8)
  {\texttt{M01\_FULL}\\共同自身分支 $+$ 逐边分支\\$\rightarrow$ 正值预测};
\draw[arrow] (common.south -| own.north) -- (own.north);
\draw[arrow] (common.south -| edge.north) -- (edge.north);
\draw[arrow] (own) -- (matched);
\draw[arrow] (edge) -- (full);
\draw[arrow] (own.south) -- ++(0,-.12) -| ([xshift=-12mm]full.north);
\node[box,fill=ThesisPeach,text width=12.8cm] at (0,-5.55)
  {统一比较：输出链接、训练损失、时间切分、评价域与失败回退\\
  底层参数对照另含 $t$ 与 $t-5$ 的归档参数及同样的有效性标志};
\end{tikzpicture}
\caption{完整模型与风险汇总基线的信息匹配结构。虚线框标记无逐边分支的基线；颜色辅助区分模型，不编码实证优劣。}
\label{fig:matched_architecture}
\end{figure}


\subsection{边条件消息传递}
将节点输入映射到$d$维表示：
\begin{equation}\label{eq:node_encoder}
h_{i,t}=\relu(W_z z_{i,t}+b_z).
\end{equation}
对来源$j$到目标$i$，计算逐维门控和消息：
\begin{align}
a_{j\to i,t}^{\mathrm{gate}}&=\sigmoid\!\left(f_{\theta}([h_{i,t},h_{j,t},e_{j\to i,t}])\right),\label{eq:gate}\\
m_{i,t}&=\frac1{\max(1,|\mathcal N_i(t)|)}
\sum_{j\in\mathcal N_i(t)}a_{j\to i,t}^{\mathrm{gate}}\odot W_mh_{j,t}.\label{eq:message}
\end{align}
门控为$d$维向量，$\odot$表示逐元素乘法。主开发版本使用一层消息传递和$d=32$。采用sigmoid而非入边softmax，避免强制归一化在某些情况下抵消共同风险强度；该选择本身也通过替代聚合规则进行检验。没有有效邻居时$m_{i,t}=0$。

式\eqref{eq:message}使边特征调节具体来源的消息，而不是先将边全部平均成一个节点数字。其一般思想属于边信息参与的消息传递框架\cite{gilmer2017}，本研究不宣称发明新的通用消息传递范式。

\subsection{自身分支、图分支与正值预测}
完整匹配版本的预测为
\begin{align}
\eta_{i,t}^{(h)}&=b_{i,h}+\beta_h^{\top}\widetilde x^{H}_{i,t}
+f_{\mathrm{own},h}(z_{i,t})+w_h^{\top}m_{i,t},\label{eq:linear_predictor}\\
\widehat y_{i,t}^{(h)}&=s_{i,c}\left[\varepsilon+\softplus(\eta_{i,t}^{(h)})\right],\quad\varepsilon>0.\label{eq:positive_prediction}
\end{align}
$f_{\mathrm{own}}$为浅层自身非线性分支。匹配风险汇总基线\code{B06_MATCHED}具有完全相同的节点输入、$f_{\mathrm{own}}$、正值链接及训练协议，仅删除$w_h^\top m_{i,t}$。由此，完整模型与汇总基线的差别可以被明确定位为逐边分支。对参数数量差异另用等容量自身MLP控制，防止将容量收益当作关系收益。

最小版本将节点信息限制为HAR，并令$f_{\mathrm{own}}=0$，保留自身线性结构与邻居分支。使用softplus的匹配模型不是与传统OLS完全相同的HAR，因此本文分别报告传统基线和统一架构对照。

\subsection{底层参数与表示增量}
给定当前及五个交易日前的$\rho$、copula自由度、两端$\sigma$和边际自由度，可以通过式\eqref{eq:conditional_t}—\eqref{eq:edgefeatures}确定性重建主边特征。因此，本文设置底层参数边模型，将这些当前及滞后量直接输入相同骨干，并保留相同有效性标志。

只给当前参数不足以重建第五日差分项，因而不属于等信息比较。主模型若优于完整参数对照，支持的是风险坐标的归纳偏置或有限样本表示效率，而不是从确定性变换中获得了新的原始信息。

\section{损失函数与估计目标}
对$y>0$、$v>0$，定义
\begin{align}
L_{\mathrm{MSE}}(y,v)&=(y-v)^2,\label{eq:mse}\\
L_{\mathrm{Q}}(y,v)&=\frac yv-\log\frac yv-1.\label{eq:qlike}
\end{align}
二者均以较小为优。训练时QLIKE可采用去除与预测无关项的形式$y/v+\log v$，但主表不将该可能为负的值与式\eqref{eq:qlike}混用。真实有效零标签另行报告，不能以停牌或缺失填零后进入正常评价。

为改善数值尺度而逐股缩放输入与输出，不应隐含改变MSE的股票权重。主MSE在原始方差单位上计算，至多除以对所有股票一致的训练全局常数。QLIKE在同时缩放$y,v$后保持不变，故可在缩放坐标上稳定计算。

需要强调，QLIKE并不把条件均值预测改成尾部分位数预测。在预测器不受函数类限制且相关期望存在时，对固定信息集的去常数风险求导，有
\begin{equation}\label{eq:qlike_optimum}
\frac{\partial}{\partial v}\E\left[\frac{y}{v}+\log v\mid\F_t\right]
=\frac{v-\E[y\mid\F_t]}{v^2}.
\end{equation}
因此最优正值$v$仍是$\E[y\mid\F_t]$。不同损失在受限模型、有限样本和优化过程中可以表现不同，不能据此把本文称为未来RV尾部分位数预测模型。

\section{本章小结}
本章明确了三个对象：未来日均已实现方差、日收益条件分布构造的风险边，以及用边调节邻居信息的正值预测器。完整模型与汇总基线共享自身历史、公共信息和风险汇总，使逐边增量接受直接比较。尾部镜像、尺度分解和底层参数对照共同限定了风险边的解释范围。

\chapter{数据说明与实证设计}
\section{数据来源与当前可证实范围}
\subsection{公司数据的研究用途}
本文计划使用公司维护的中国股票分钟行情与配套日频、证券状态数据。分钟价格用于构造日度已实现方差，日收益用于估计条件风险关系，两者在相同证券与交易日历上对齐。主输入优先采用一分钟价格序列统一构造五分钟测度，存量五分钟数据仅在独立认证后作为对照，不以不同来源的历史数据静默拼补。

本版已读取原始数据位置索引及五份数据字典并记录文件哈希；本轮登记对象的实际字段结构均已取得，并已通过授权只读查询获取有界真实小样，公司行动严格版本视图也已完成有界读取。小样获取成功不等于验收通过：初步核对已识别收益来源口径差异及历史时点可见性证据不足，分钟会话规则仍在核验，逐股逐日覆盖尚未确认。字段存在及有限记录可读不能证明完整历史可用。表\ref{tab:data_status}说明后续验收事项；正式协议尚未锁定，具体日期、证券数量和覆盖统计仍待确认后填入。

\begin{table}[htbp]\centering\small
\caption{数据来源类别与验收状态}\label{tab:data_status}
\begin{tabularx}{\textwidth}{p{0.2\textwidth}YY}
\toprule 数据类别 & 已报告的用途与性质 & 实证前待确认事项\\\midrule
股票分钟价格 & L2成交价格快照派生，OHLC与区间量额 & 真实覆盖、区间标签、无成交规则、端点价格\\
股票日线 & 日收益和量额核对，含交易所参考前收 & 时间编码、参考价与公司行动一致性\\
公司行动 & 含市场可见时间的版本化视图，已取有界小样 & 历史覆盖、事件与条款版本联结、可见性证据\\
交易日历与状态 & 日历、停复牌、限价及质量记录 & 有效水位、缺口与哨兵值含义\\
历史成分与行业 & 带来源的成分、行业有效区间 & 发布证据、分类版本、历史修订\\
证券身份 & 代码、市场、上市与退市字段 & 当前快照限制与跨代码主体歧义\\\bottomrule
\end{tabularx}
\end{table}

本研究不需要客户账户、公司交易策略、持仓或逐笔委托数据。内部连接细节由独立的数据适配文档管理，论文仅说明可披露的数据类别、频率、时段与处理方法。最终是否公开派生RV和风险边，也须获得明确的数据授权。

\subsection{三个不同层次的时点保证}
本文将时点问题分为三个层次。第一，算法时间顺序，要求起点$t$的特征、图和预测不能读取$t$之后的观测。第二，事件及分类可见性，要求使用可证明在当时公开的版本，而不是将后来的修订提前用于历史。第三，行情本身的历史版本与实际入库时刻，即当前历史库是否还原了当时实际可读的记录。

本稿计划对第一层实施严格测试，并在公司行动等具备证据的部分使用第二层。第三层目前尚未取得完整证明。因此，即使所有算法都按时间顺序运行，也只能在相应数据快照上声称时间外预测，不能自动声称完整复原了历史线上可得信息。程序运行时间不是历史发布时点，后补的历史记录也不能仅凭业务日期获得当时可见的资格。

原始字典声明，交易日$D$的分钟数据最早在$D+1$日北京时间05:00之后、日线最早在06:10之后，且相应任务成功后才可用。这是待实测确认的文档条件，不是已验证的服务时效保证，也不能据此称完整输入在$D$日收盘时即可读取。第三章中的$t$首先表示最后纳入观测的交易日；真正发出预测的时刻还须晚于全部输入的可用时刻。正式协议须确认该时刻是否早于目标首日的交易会话；若不能满足，需在测试前调整预测起点与标签对齐，而不能把延迟到达的数据提前使用。

\section{样本选择与历史覆盖}
\subsection{样本形成规则}
开题提出约30只具有流动性和行业代表性的沪深300股票作为主样本。本文保留这一规模目标，但将形成日设在首个预测模型训练起点之前：先用历史预热段确定可估计风险关系的候选证券，在形成日根据当时可证明的成分资格、此前成交额、历史可用性和行业信息选样，然后固定集合开展后续训练、验证和测试。

这种安排不使用测试期是否持续交易、是否退市或模型预测效果进行反向选择。样本形成后发生的停牌、成分退出和证券身份变化均作为后续数据状态记录，不因“不方便”而将该股票从整个历史中删除。形成日前的预热数据用于选择和初始化，不作为已经对该固定样本开展在线预测的证据。

如果历史成分的公布时间不可证实，则研究应降级为明确限定的固定股票样本，不再声称严格的历史沪深300可投资成分筛选。若完整跨代码证券身份不可得，则采用交易所、代码与上市身份区间构造临时键，禁止仅凭股票名称串接不同主体。

\subsection{历史行业与共同信息}
行业匹配使用有效区间和分类版本共同确定。名称含某一版本年份的表并不代表全部历史都采用该版本，追溯性分类也不自动等于当年已知信息。不能以当前行业成员替代历史成员，并将由此形成的行业平均值当作完全无前视的控制。

在历史行业无法认证时，匹配组统一删去行业变量，但仍可使用剔除目标股票的样本平均RV：
\begin{equation}\label{eq:market_control}
M_{-i,t}=\frac{1}{|\mathcal V_t\setminus\{i\}|}
\sum_{j\in\mathcal V_t\setminus\{i\}}\RV_{j,t},
\end{equation}
其中$\mathcal V_t$为当日测度有效的股票集合，并同时记录有效数量。该变量是样本中个股RV的平均值，不是市场指数RV；后者还受到成分协方差及权重影响。市场指数分钟数据尚未认证，不阻塞采用明确定义的共同信息对照。

\subsection{覆盖统计与样本长度}
较长历史用于容纳风险网络初始化、预测模型训练、验证和测试。初始预算候选为750个交易日的风险预热、750个训练起点、250个验证起点和500个测试起点。它们是开发规划，不是已查得的覆盖或通用最低样本定理；实际还要扣除缺失、特征回看和标签成熟造成的损失。

覆盖统计至少包含原始分钟记录、有效五分钟端点、完整股票日、连续HAR窗口及一、五、二十二日成熟标签。仅报告最早与最晚日期，不足以证明中间没有缺口；同日30只股票也不能直接视作30个独立市场时点。若真实历史不足，须在最终测试之前调整规模或切分，并明确统计精度的限制。

\begin{table}[htbp]\centering\small
\caption{正式样本待验收信息}\label{tab:sample_pending}
\begin{tabularx}{\textwidth}{p{0.34\textwidth}Y}
\toprule 项目 & 本版状态\\\midrule
样本股票与证券身份列表 & \Pending{实际名单和形成日}\\
分钟原始区间与有效交易日 & \Pending{逐股逐日覆盖统计}\\
收益调整证据与公司行动覆盖 & \Pending{核对结果和版本限制}\\
训练、验证与测试区间 & \Pending{冻结日期与数据版本}\\
各期限有效标签及排除原因 & \Pending{样本损失分解}\\\bottomrule
\end{tabularx}
\end{table}

\section{数据预处理与测度验收}
\subsection{时间编码与会话边界}
已报告的源合同存在重要差异：一分钟导出记录的UTC样式后缀可能只是裸北京时间的格式标记，而旧五分钟源的同一后缀才表示真实UTC。因此，不能对所有表统一使用“转换到北京时间”的逻辑。

实际接入按来源与传输接口分别建立时间profile，比较原始字符串或epoch、数据库精度、会话时区和解码后的开盘、午休及收盘位置。一旦内部统一为带时区的北京时间，禁止再次解释导致八小时偏移。日线与指数数据也单独认证，不直接继承分钟表的时间规则。

五分钟价格路径在上午和下午分别建立。开盘集合竞价单独成根时，应明确其作为锚点还是进入区间汇总，并避免重复计算；下午首价的真实时点需记录。价格保持仅限于已证实的无成交区间，不能用后续价格反填缺失，也不能跨休市期间生成不存在的成交区间。

字典另声明一分钟记录采用右端标签、09:30开盘集合竞价单列、15:00记录包含收盘集合竞价。这些声明需要用授权小样及真实schema核对后才进入有效采样规则；本稿不将它们当作已经观察到的数据库事实。

\subsection{价格、量额与异常状态}
数据检查包括同键重复、冲突版本、价格正值、OHLC约束、量额单位和交易状态。完全相同的重复可以按固定规则去重；同键不同值的冲突需保留证据，不按文件最后一行直接覆盖。区间股数与区间成交额先与日线核对口径，再评价差值。

真实极端波动、涨跌停及高成交额本身不是数据错误。本文主标签不按全样本分位数缩尾，也不为降低MSE删除真实高波动日。停牌、无成交和数据缺口分开处理：停牌不自动等于零方差；没有记录不自动证明没有成交；零成交量下重复的最新成交价是否可靠，需由源生成逻辑和样本确认。

\subsection{收益调整与公司行动}
候选参考价收益式\eqref{eq:reference_return}需对普通日、除权除息日及其他特殊事件进行核验。若采用价格调整因子$A_{i,t}$，定义
\begin{equation}\label{eq:adjusted_return}
R^{\mathrm{adj}}_{i,t}=\log\frac{P_{i,t}A_{i,t}}{P_{i,t-1}A_{i,t-1}}.
\end{equation}
它与参考价收益一致所需的因子方向为
\begin{equation}\label{eq:factor_direction}
\frac{A_{i,t}}{A_{i,t-1}}=\frac{P_{i,t-1}}{B_{i,t}}.
\end{equation}
数据库字段名为ratio并不保证恰好定义为式\eqref{eq:factor_direction}，必须核对实际说明与样例。现金分红、股份变动和权利到账还会影响财富回报，因此价格调整恒等式不构成总回报认证。

版本化公司行动需同时满足两个时间条件：使用在预测截止之前已知的事件版本，并在经济生效时处理相应价格机械变化。不能在公告当天提前实施未来除权调整，也不能只获取最新事件快照回填过去。事件覆盖表用于区别“确实无事件”和“尚未覆盖”，没有查到事件行不能自动得出前者。

\subsection{数据质量与模型有效样本}
每个股票日保存测度有效性、特征有效性和标签有效性。特征缺失不应自动删除其他股票的当日记录；图边失效也不能成为主模型独占的删样理由。所有模型首先在同一标签域内生成预测，对失败使用预设兜底，并分别报告失败率和兜底占比。

在实现验收中，应选取真实完整股票日手算RV，并核对至少一种特殊会话或公司行动情形。纯构造数据用于测试公式和边界，但不代替真实数据清洗验证，更不进入主市场结果表。

\section{实证模型与贡献归因}
\subsection{模型体系}
\begin{table}[htbp]\centering\small
\caption{主比较模型及其作用}\label{tab:models}
\begin{tabularx}{\textwidth}{p{0.23\textwidth}YY}
\toprule 模型组 & 主要信息/结构 & 对应比较目的\\\midrule
传统HAR、HARQ & 自身多尺度历史、四次变差 & 强计量基准\\
匹配正值HAR、自身MLP & 相同缩放、链接与优化 & 损失和非线性容量\\
公共信息增强模型 & 市场/可用行业/边际风险 & 排除共同风险解释\\
风险汇总MLP & 相同公共信息和风险汇总 & 核心关系增量基线\\
Ridge-VHAR & 正则化横截面HAR输入 & 非图跨股票竞争者\\
GHAR、GLASSO-GNN & 线性/非线性统计图 & 图骨干与一般关联\\
DCC相关边模型 & 相同动态相关来源 & 排除仅有动态性收益\\
CoVaR二值/标量模型 & 固定支持、不同边输入 & 拓扑与边向量区分\\
底层参数边模型 & 当前及滞后原始参数 & 风险表征归纳偏置\\
完整边特征模型 & 相同自身/汇总+逐边消息 & 主确认性模型\\\bottomrule
\end{tabularx}
\end{table}
传统基线用于文献比较，统一架构组用于内部归因。GLASSO原文复现与统一骨干变体单独注册；忠实复现不能因使用相近名字就省略输入区间、归一化和损失的核对。相关构图使用何种收益、是否保留负关联，以及如何匹配边预算，均写入配置。

\subsection{主实验与消融顺序}
实证依次开展四个层次。第一层比较自身模型在MSE与QLIKE训练下的表现。第二层依次加入公共信息、风险汇总和逐边分支，检验完整模型是否超过匹配风险汇总基线。第三层保持节点信息、连接预算和骨干一致，替换动态相关、CoVaR拓扑、标量边与向量边。第四层通过尺度标准化、底层参数和关系破坏实验，检验改善是否依赖特定风险语义和邻居对应关系。

关系破坏包括固定拓扑重配入边特征、保度或近似保度重连、对称化及使用滞后图。作为正式消融时重新训练全部模型；仅在推理时改变边的结果则作为敏感性分析，不能与重训消融混合。随机重连改变了哪些结构、转置是否改变入度，均需说明。

动态性另拆为依赖状态更新与边际尺度更新：固定其中一部分、仅更新另一部分，再与全部更新比较。这样避免将“波动尺度在变”全部解释为“依赖网络在变”。上行通道若由数学性质决定重复，就不进行两个独立信号的贡献宣称。

\section{样本外训练与预测协议}
\subsection{时间切分与标签成熟}
设$c$为一次模型训练截止日，允许进入拟合的样本满足
\begin{equation}\label{eq:training_set}
\mathcal T_c^{(h)}=\{(i,s):\ \mathrm{end}(s,h)\leq c,
\ \mathrm{available}(y_{i,s}^{(h)})\leq c,\ \mathrm{valid}_{i,s,h}=1\}.
\end{equation}
这里的available只有在具备真实可得时间证据时才直接使用；否则明确采用固定历史快照下的成熟规则，并披露不能验证入库延迟。仅起点$s$已经过去，并不表示未来22日标签已经成熟。

训练、验证和最终测试按连续时间切分，不能随机打散股票日或图快照。跨边界目标窗口按最大期限22日的共同规则隔离，另可报告期限专属规则。每次拟合保存训练起点范围、目标终点范围和实际可用日期，使标签泄漏可以被自动检测。

\subsection{滚动更新与调参}
开发段为每个学习模型族提供至多12组候选设置，以两个固定随机种子选择结构与训练轮数。默认开发起点为一层图、隐藏维度32、学习率$10^{-3}$、权重衰减$10^{-4}$和dropout 0.1，均是待验证的工程参数，而非已观察到的最优值。

正式测试约每21个交易日重训，使用开发期确定的训练轮数，并按同一规则纳入已成熟的过去样本。不得根据最终测试表现修改图密度、删除年份、改变主要期限或重选胜出的随机种子。五个固定种子用于反映优化随机性，确定性基线不通过复制运行变成更多独立市场证据。

\subsection{历史不变性与失败规则}
将$t$之后的行情修改或追加后，在同一已冻结证券集合和随机状态下，截至$t$的特征、图和预测应保持不变。该测试同时覆盖价格处理、标准化、构图和训练数据索引。当前时点重估参数可以改善当前状态，但不能覆盖更早起点已归档的图。

图估计失败时优先使用年龄未超限的合法参数重新过滤；不可用则缺边。预测层保留无图回退。若无图模型也不可用，使用预设的近期有效RV或训练历史均值，并记录原因。主评价保留这类操作后预测，防止模型通过不输出困难样本改善结果；纯成功子集仅作附表。

\section{评价指标与统计检验}
\subsection{日期级损失与相对改善}
对模型$m$，设$I_t^{(h)}$为所有模型共同的目标有效股票集合，$S$为随机种子数，定义日期级模型族损失
\begin{equation}\label{eq:daily_loss}
\ell_{m,t}^{(h)}=\frac1S\sum_{s=1}^{S}\frac1{|I_t^{(h)}|}
\sum_{i\in I_t^{(h)}}L\big(y_{i,t}^{(h)},\widehat y_{m,s,i,t}^{(h)}\big).
\end{equation}
首先对同日截面与种子取平均，再以日期为统计单位。种子平均损失不等于先平均预测再计算损失；后者属于另一个集成模型，需要单独登记。

记基线为$b$、主模型为$m$，定义
\begin{equation}\label{eq:difference}
d_t=\ell_{b,t}-\ell_{m,t},\qquad
\Delta_{b,m}=\frac1T\sum_t d_t,\qquad
\mathrm{Gain}_{b,m}=1-\frac{\bar\ell_m}{\bar\ell_b}.
\end{equation}
正的$\Delta$或Gain表示主模型损失较低，但Gain是损失的相对降低，不是收益率或交易利润。主指标为QLIKE，MSE、RMSE和MAE辅助报告，所有表明确单位和有效样本数。

\subsection{区块置信区间与 DM 检验}
日期级差值保留了同日股票共同冲击。主置信区间采用成对日期区块重采样，同时保留不同模型在同一天的对应关系。候选设置为2000次重采样、平均块长50个交易日，并以25和100作敏感性；这些设置在测试之前冻结。

DM比较以损失差均值为对象\cite{diebold1995}。采用HAC长期方差估计$\widehat\Omega$，统计量为
\begin{equation}\label{eq:dm}
\mathrm{DM}=\frac{\bar d}{\sqrt{\widehat\Omega/T}},\qquad
\widehat\Omega=\widehat\gamma_0+2\sum_{k=1}^{K}\left(1-\frac{k}{K+1}\right)\widehat\gamma_k.
\end{equation}
候选带宽$K=\max(h-1,20)$并作敏感性分析。多日目标重叠、滚动参数估计和厚尾损失差会影响有限样本推断，因此报告效应量、置信区间及适用限制，不只报告一个$p$值。

\subsection{多重比较、MCS 与市场状态}
主比较预设为完整模型对匹配风险汇总、公共信息、动态相关、标量CoVaR和底层参数模型等少量差值，并对比较族实施Holm校正。MCS在相同日期损失矩阵上计算90\%与95\%集合\cite{hansen2011}。进入集合并不表示显著优于集合内其他模型。

高风险状态由起点$t$已经可得的市场汇总RV与训练确定的阈值比较得到，例如开发历史的90\%分位数。未来实际极端日可以用于事后误差剖析，但不进入实时状态变量。高、低状态差异应直接检验差值之差，不能以“高状态显著、低状态不显著”替代组间差异检验。

\section{稳健性、计算成本与本章小结}
预设稳健性覆盖风险估计窗口、图稀疏度、采样频率、点条件与事件条件、边际模型及样本规模。所有参数不做无目的全笛卡尔积，而围绕核心结果分层执行。图估计对所有模型和种子共享缓存；报告边际与copula估计时间、模型训练时间、推理时间、失败率和内存需求，以评价复杂度与预测收益之间的取舍。

本章将数据位置转化为可验收的数据协议，并将研究问题转化为有明确对照、时间范围和统计单位的实验。下一步先验证真实分钟形态与收益口径，再锁定样本和切分。只有完成这些环节，第五章中的实证数字才具有明确的数据与方法含义。

\chapter{实证结果与分析}
\noindent\textbf{本章状态：}本版尚未运行公司数据的正式市场实验。以下内容给出结果组织、表格位置和对应论证问题；破折号表示未生成结果，不表示零损失或无显著差异。正式运行后，应以结果文件替换占位内容，并据真实证据完成段落。

\section{样本与风险网络的描述性结果}
\subsection{数据覆盖及主要变量}
本节首先报告最终样本形成日、证券构成、有效区间和各处理阶段的样本数量，使读者能够看清从原始分钟记录到未来方差标签的变化。描述性统计包括RV、参考价收益、条件sigma、CoVaR强度和有效边比例，并将缺失、停牌、未成熟标签及模型失败分别列示。

\Pending{表5.1：逐阶段样本损失与变量统计；图5.1：逐股逐日覆盖；图5.2：训练、验证、测试及成熟边界。}

风险网络描述先比较原始和标准化强度，再讨论条件输入、输出差异。网络图只辅助展示若干预先确定状态，固定节点位置与边尺度；不能用布局变化制造结构变化的视觉印象。

\subsection{风险定义与代数检查}
本节列示上下尾镜像误差、同形边际下的标准化双向差、不同边际自由度的分解及图支持非对称率。这些诊断用于确认实现符合第三章的数学设定，不把标准t-Copula的结构对称性当作数据中“发现没有非对称风险”的结论。

\Pending{表5.2：数值与边语义诊断；图5.3：原始方向差与标准化方向差。}

\section{整体样本外预测表现}
表\ref{tab:main_results_pending}保留全部主要模型与期限。正文应首先报告相对于强基线的差值和区间，再说明最优均值排序是否稳定；不能只展示表现最好的种子或删去主模型没有获胜的期限。

\begin{table}[htbp]\centering\small
\caption{主预测结果表结构：QLIKE（尚未生成实证值）}\label{tab:main_results_pending}
\begin{tabularx}{\textwidth}{Yccc}
\toprule 模型 & 1日 & 5日 & 22日\\\midrule
传统HAR / HARQ & — & — & —\\
匹配自身MLP & — & — & —\\
公共信息增强模型 & — & — & —\\
匹配风险汇总MLP & — & — & —\\
正则化横截面HAR & — & — & —\\
GLASSO图模型 & — & — & —\\
DCC相关边模型 & — & — & —\\
CoVaR标量边模型 & — & — & —\\
底层参数边模型 & — & — & —\\
完整边特征模型 & — & — & —\\\bottomrule
\end{tabularx}
\end{table}

正式表格同时标注训练损失、节点信息组、有效股票日及日期数。MSE、RMSE、MAE另表给出，不把不同单位的改善画在同一数值轴上。可用相对损失便于阅读，但仍保留原始损失及相应分母。

\section{预测改善来自何处}
\subsection{训练损失与公共信息}
首先比较同一模型从MSE训练到QLIKE训练的变化，再比较同一QLIKE训练下无图、公共信息增强和风险汇总模型。这样可判断总体改善中有多少在使用具体图结构之前已经出现。

\Pending{表5.4：训练损失交叉比较；图5.4：损失准则与公共信息贡献拆分。}

若完整模型仅超过MSE训练的HAR，但未超过同损失强基线，结论应限于整体流程或训练方式的改善，不将全部差异记为逐边信息贡献。

\subsection{风险汇总与逐边条件化}
核心比较为\code{M01_FULL}与\code{B06_MATCHED}。两者共享节点信息、风险汇总、自身MLP和训练方案，差别在逐边分支。应报告平均损失差、相对降低、日期区块置信区间以及不同年份的稳定性，并检查预测失败或缺边是否改变有效域。

\Pending{表5.5：核心成对比较与多重校正；图5.5：相对匹配汇总基线的累计日期损失差。}

若逐边分支的改进得到支持，说明具体关系在研究样本中保留了汇总变量之外可用于预测的结构；若区间较宽，则只能说明目前证据尚不足以精确区分，而不将点估计排序当作稳定优越性。

\subsection{拓扑、标量与多维边}
在固定支持下比较二值边、标量边和向量边，进而删除标准化、反向和变化通道。对每个消融报告重训还是推理屏蔽、参数量与入度是否变化。关系重配与随机重连结果用于检查“边值属于哪个邻居”是否重要，不将这些计算实验称为真实市场因果干预。

\Pending{表5.6：边特征与关系破坏消融；图5.6：成对差值及区间。}

\subsection{边际尺度与风险表征}
比较标准化风险边、动态相关边以及包含当前和滞后底层参数的模型。若风险表征优于完整参数对照，讨论其有限样本表示效率；若收益被sigma或公共风险解释，则按相应层次限定贡献。上下尾或反向通道的代数重复，不应通过重复检验获得多项独立贡献的表象。

\section{期限、市场状态与稳健性}
本节按预设的1、5、22日目标和预测时可知状态报告结果。状态表需包含实际日期数量和成对差异的组间检验。年份分析与累计损失差共同用于检查结果是否由少数日期驱动，不只展示模型最有利的危机区间。

稳健性表覆盖5/10/15分钟采样、500/750/1000日风险窗口、不同入边数及点/事件条件定义。每项结果附实际有效样本；若改变测度或窗口导致评价域变化，应再提供共同域比较。计算成本单独报告，不将更多算力投入隐藏在方法名称之后。

\Pending{表5.7：期限与状态；表5.8：稳健性与成本。}

\section{本章小结的回填原则}
正式小结应依次回答：是否超越强自身基线、是否超越公共信息和风险汇总、哪些边内容有独立增量、效果在哪些期限与状态成立、复杂度代价是多少。每项结论对应明确结果表，不以方法名称或图示替代证据。本稿不预先填写这些市场结论。

\chapter{研究结论与实践启示}
\section{研究框架与待验证结论}
本文围绕中国股票已实现波动率预测，提出由高频测度、动态条件风险关系与边特征消息传递构成的研究框架。与仅比较复杂模型和简单HAR不同，本文将风险汇总、具体关系、动态相关和边际尺度放入统一的贡献归因协议，以检验逐边条件化是否具有额外价值。

在方法层面，本稿已经明确三个结论边界。首先，所选对称边际与t-Copula设定中的上下尾强度具有镜像对应，不能预先作为两个独立信号。其次，原始条件风险量的方向差可能包含接收方波动尺度的作用，标准化与分布形状控制是解释方向的重要步骤。最后，风险汇总在一般图结构中不能决定逐边加权邻居状态，但这种表示差异是否有实际预测收益，必须由有限样本实验确认。

上述内容属于模型定义与数学推导，不替代市场实证。最终结论应在正式结果生成后分别说明：完整模型相对匹配风险汇总与动态相关模型的增量、有效边特征、适用期限和状态，以及计算与数据条件。

\Pending{根据第五章核心表格回填主要实证结论；每条结论注明适用样本、效应大小和不确定性。}

\section{实践启示}
\subsection{风险预测与数据治理}
对实际使用而言，风险预测结果应与数据时段、收益调整口径和失败标志一起解释。只给出模型输出而不说明其输入是否完整，容易把数据缺失造成的平稳误判为真实低风险。公司数据适配中的时间编码、参考价、公司行动版本与历史成分证据，因而不是模型之外可忽略的技术细节。

若逐边模型经样本外比较显示稳定增量，其适当用途是为既有风险预测增加条件关系信息，而不是把高权重边直接命名为市场中的因果传染路径。若风险汇总已经取得相近效果，则较简约的模型在运行成本和维护方面可能更合适；这一取舍应结合实际差值与成本报告判断。

\subsection{资产配置及监管解释的边界}
本文的主目标是方差预测，而非交易收益最大化。将方差预测用于组合权重、VaR或ES，还需额外的收益分布、协方差结构、持有区间及执行约束。没有对应回测证据时，不将预测误差降低写成已经提高净收益、降低回撤或识别系统重要性主体。

对于市场监测，条件风险网络可以作为描述和研究工具，但其解释应保留模型假设、样本代表性和边际尺度限制。固定股票样本中的输入、输出排名不能直接外推为整个市场的风险贡献排序。

\section{研究局限与未来方向}
本研究的首要数据局限是公司历史分钟与辅助数据的版本可得性。若只有当前修订的历史快照，严格的算法时间顺序仍不能完全还原历史线上环境；这一限制需在最终论文中保留。其次，二元关系模型的拼接不保证全维联合分布有效，因而本研究不直接解决完整协方差与组合风险分配问题。

模型方面，标准t-Copula限制了上下尾非对称的表达，浅层图模型则在表示能力与有限样本稳定性之间作出取舍。未来可以在保持同一强基线和时间协议的前提下，比较非对称依赖、更多证券以及不同交易时段的外部效度，而不是仅因模型更复杂就视为自然改进。

在最终研究结论中，应把方法结构所保证的性质、真实样本所支持的结果和仍需验证的应用分别陈述。这样的区分使论文从研究问题到证据形成闭合，也为后续复现和扩展保留清楚的起点。

\printbibliography[heading=bibintoc,title={参考文献}]
\appendix
\chapter{补充推导与实现对应}
\section{对称设定下的上下尾镜像}
本节证明命题\ref{prop:symmetry}。在本文的连续t-Copula中，联合密度存在且分位数唯一。设目标概率坐标为$V$，来源为$U$，其copula密度满足中心对称性
\begin{equation}
c(u,v)=c(1-u,1-v).
\end{equation}
由于$U$的边际密度为1，$V$在$U=\alpha$下的条件分布函数为$H_{\alpha}(v)=\int_0^v c(\alpha,w)\,dw$。作变量代换得到
\begin{align}
H_{1-\alpha}(1-v)
&=\int_0^{1-v}c(1-\alpha,w)\,dw\\
&=\int_v^1c(\alpha,z)\,dz
=1-H_{\alpha}(v).
\end{align}
因此$H_{1-\alpha}^{-1}(1-\beta)=1-H_{\alpha}^{-1}(\beta)$。目标边际关于$m_i$对称意味着$F_i^{-1}(1-v)=2m_i-F_i^{-1}(v)$，故
\begin{equation}
q_{i\leftarrow j}(1-\beta\mid1-\alpha)=2m_i-q_{i\leftarrow j}(\beta\mid\alpha).
\end{equation}
同理，当参考态为$\alpha=0.5$时，该关系依然成立。将困境态与参考态相减，得到$\Delta^+=-\Delta^-$。这一性质不要求两只股票具有相同边际自由度，只要求各自用于该条件方向的目标边际对称。上下行都使用同一预测起点的参数，因此参数随时间变化不会打破逐起点的镜像性质。

若上、下尾阈值不互为镜像，或两侧使用不同参考态，则不能套用这个等式。数值实现中的差异应先检查这些定义，而不是直接归因于市场非对称性。本稿的参考测试覆盖不同相关系数、边际自由度、条件位置和尺度；其用途是检查实现与命题的一致性，而不是替代真实市场检验。

\section{条件方向的尺度分解}
设两端条件边际为$F_i(x)=G((x-m_i)/\sigma_i)$和$F_j(x)=G((x-m_j)/\sigma_j)$，其中$G$相同。交换来源与目标时，由copula可交换性，相同$\alpha,\beta$得到同一条件概率位置$v_{\alpha,\beta}$，中位参考态位置为$v_{0.5,\beta}$。于是
\begin{align}
\Delta^-_{i\leftarrow j}
&=\sigma_i\left\{G^{-1}(v_{\alpha,\beta})-G^{-1}(v_{0.5,\beta})\right\},\\
\Delta^-_{j\leftarrow i}
&=\sigma_j\left\{G^{-1}(v_{\alpha,\beta})-G^{-1}(v_{0.5,\beta})\right\}.
\end{align}
令$k_{ij}=G^{-1}(v_{0.5,\beta})-G^{-1}(v_{\alpha,\beta})$，即得命题\ref{prop:scale}。位置参数在增量中抵消。这个抵消也说明，主边特征的底层参数控制不需要为了重建增量而额外引入条件均值；不过若改为使用CoVaR水平本身，则条件均值必须加入对应控制。

当同一目标的正强度满足$s_{i\leftarrow j}=\sigma_i k_{ij}$时，行和归一化权重为
\begin{equation}
\frac{s_{i\leftarrow j}}{\sum_{\ell}s_{i\leftarrow\ell}}
=\frac{k_{ij}}{\sum_{\ell}k_{i\ell}},
\end{equation}
因此不能同时声称“该归一化去除了个股尺度影响”与“归一化后的边仍保留了全部原始尺度信息”。类似地，top-$k$所选关系可能不互相保留；因此即使$k_{ij}=k_{ji}$，支持矩阵$A$仍可能不对称。正文的方向分析同时报告权重与支持，避免混淆。

若两端标准化边际形状不同，目标逆分位变换也不同，尺度标准化不能保证双向完全相等。本稿要求保留两端边际自由度，并进行共同边际形状控制，而不是将标准化后残留的一切差异都解释成独立方向。

\section{汇总不决定逐边条件信息：一个构造}
以下数字仅为代数示例，不是公司数据或实证结果。设节点1、2为来源，节点3、4为目标，唯一四条边的权重为
\begin{equation}
W=\begin{pmatrix} a&b\\ b&a\end{pmatrix},\qquad
W'=\begin{pmatrix} b&a\\ a&b\end{pmatrix},
\end{equation}
其中行对应目标3、4，列对应来源1、2。两个网络的每个目标入权重和、每个来源出权重和、度数及各目标权重集合均相同。然而，当来源状态为$x=(x_1,x_2)^\top$时，两个网络对目标3的线性消息差为$(a-b)(x_1-x_2)$。例如$a=0.8,b=0.2,x_1=1,x_2=4$，该目标消息分别为1.6和3.4。

因此，即使完整保留每个节点的风险总量，也可能无法恢复哪条风险关系对应哪个来源状态。该构造仅说明一般关系表示中的信息压缩，不说明这四条边能由任意给定的t-Copula参数生成，也不证明其预测价值。在真实实验中必须通过同信息集的汇总模型、等容量模型及关系重配实验验证。

\section{参考价收益、调整收益与总回报}
令$P_{t-1}$为前一交易日未复权收盘，$P_t$为当前收盘，$B_t$为当日交易所参考价。若复权价格定义为$\widetilde P_t=A_tP_t$，则其对数收益为
\begin{equation}
R_t^{\mathrm{adj}}=\log\frac{P_t}{P_{t-1}}+\log\frac{A_t}{A_{t-1}}.
\end{equation}
它等于候选参考价收益$\log(P_t/B_t)$的条件为
\begin{equation}
\frac{A_t}{A_{t-1}}=\frac{P_{t-1}}{B_t}.
\end{equation}
数据库中名为ratio的字段可能保存此比值或其倒数，累计因子也可能采用不同的锚定方向；必须从定义与样例确认，不能只凭名称代入。

仍用纯构造例子说明：若前收10元、每股现金分红1元、参考价9元且当日收盘9.9元，则参考价对数收益为$\log(9.9/9)$，若暂忽略税费和再投资，持有股票并收到1元现金的单期财富对数增长为$\log((9.9+1)/10)$。二者分别对应10\%与9\%的简单变化，因而并不相同。这不评价哪种输入对预测更好，只说明“价格调整收益”与“投资者总回报”不能未经定义直接等同。复杂送转、配股或多阶段事件更应使用明确的经济条款与生效日期。

同一会话内若所有价格统一乘以相同正因子，式\eqref{eq:intradayreturn}的差分与RV保持不变。这种不变性并不证明跨日收益正确，也不能认证整个历史复权账本的当时可得性。

\section{信息时点、标签成熟与历史归档}
设预测起点$s$对应目标区间$[s+1,s+h]$。训练截止为$c$时，标签可用至少要求$s+h\leq c$，并进一步要求目标数据已经按合同实际可用。以交易日序号表示，当$h=5,c=104$时，起点99的目标100至104已完整；起点100的目标101至105尚未成熟。判断不能依据起点100是否早于104。

历史图$G_s$由起点$s$可用的参数、状态与数据版本生成。当前$c$时点允许用已取得的历史数据重新估计当前状态，但不能将这组较新参数写回$G_s$，然后把回写后的图当成当时的输入。五日边变化必须比较$G_t$与实际归档的$G_{t-5}$；底层参数控制也从这两个归档版本取值。

证券集合在最初形成期结束时固定，早于第一条预测模型训练样本的起点。该选择只依据形成期已有的资格、流动性和覆盖信息，不根据后续完整存活或测试结果反选。若历史资格只有事后回填，数据论文必须将样本称谓与时点限制说明清楚，不能依靠“历史日期”三个字替代当时可见性证据。

\section{公式与代码的对应要求}
\begin{longtable}{p{0.30\textwidth}p{0.62\textwidth}}
\toprule 正文对象 & Codex实现及验收重点\\\midrule\endhead
式\eqref{eq:rv}、\eqref{eq:rq} & 会话锚点、完整规则采样、午休与集合竞价口径；现有参考函数只接受已整理的价格路径。\\
式\eqref{eq:target} & 用完整交易日历构造未来窗口；缺失一日不压缩窗口；保存成熟日期。\\
式\eqref{eq:unit_t} & 明确普通t与单位方差t转换，边际自由度大于2。\\
式\eqref{eq:dcc} & 正定约束、分段初始化、过去估计窗口；完整DCC估计器仍需实现。\\
式\eqref{eq:conditional_t}—\eqref{eq:quantile_back} & 从同一无序对生成双向条件量；参考测试核对镜像与尺度性质。\\
式\eqref{eq:edgefeatures} & 历史归档五日差、缺边掩码、训练尺度；底层参数对照包含同样滞后。\\
式\eqref{eq:message} & 行为目标、列为来源；编号置换等变，空邻居退回自身分支。\\
式\eqref{eq:qlike} & 正标签域、预测下限、原始与去常数损失区分；全模型共同评价样本。\\
\bottomrule
\end{longtable}
当前附带的数值代码只验证若干定义、代数性质和构造例子，未实现完整数据库读取、GARCH/DCC估计、图神经网络训练或统计实验。真实数据的验收与全部预测实验由后续研究阶段完成。

\end{document}
